\section{Análisis de resultados}

En esta sección se presentan y discuten los resultados, comparando la SNN original (A--B), la SNN híbrida (A--B--C) y modelos TSFEDL\footnote{\url{https://github.com/JGS9515/compare_to_TSFEDL}}. Se incluyen análisis por dataset, ablasiones y eficiencia.  Los tiempos de entrenamiento reportados deben interpretarse considerando que todos los modelos se ejecutaron en CPU (véase la Sección~\ref{subsec:entorno}).

\subsection{Resultados por dataset}

\subsubsection{IOPS}
Los resultados en el dataset IOPS muestran el desafío de la detección de anomalías en datos altamente desbalanceados (1.92\% de anomalías):
\begin{itemize}
    \item \textbf{Resumen}: La arquitectura híbrida SNN (A--B--C) muestra mejoras significativas sobre el modelo base en IOPS. El mejor resultado híbrido alcanza F1 = 0.277 (configuración n\_100), representando una mejora de 4.6x sobre el modelo SNN base (F1 = 0.061), aunque sigue siendo inferior a los baselines TSFEDL.
    \item \textbf{Tendencias}: Los mejores resultados se obtuvieron con kernels \texttt{mexican\_hat} y \texttt{laplacian}, procesamiento \texttt{weighted\_sum} y \texttt{max}, y configuraciones de menor número de neuronas (n\_100 > n\_200 > n\_400), sugiriendo que redes más grandes tienden a sobreajustarse en este dataset desbalanceado.
    \item \textbf{Errores}: A pesar de las mejoras, persiste el desafío del extremo desbalanceo, con el modelo híbrido logrando detectar una fracción mayor de anomalías pero manteniendo altas tasas de falsos negativos comparado con métodos tradicionales.
\end{itemize}

\begin{table}[htbp]
\centering
\small
\begin{tabular}{lccccc}
\hline\hline
\textbf{Modelo} & \textbf{N} & \textbf{Prec} & \textbf{Rec} & \textbf{F1} & \textbf{MSE} \\
\hline
SNN (A--B) & 100 & -- & -- & 0.060 & 0.142 \\
SNN (A--B) & 200 & -- & -- & 0.061 & 0.142 \\
SNN (A--B) & 400 & -- & -- & 0.062 & 0.142 \\
\hline
SNN (A--B--C) & 100 & -- & -- & 0.277 & 0.142 \\
SNN (A--B--C) & 200 & -- & -- & 0.197 & 0.142 \\
SNN (A--B--C) & 400 & -- & -- & 0.096 & 0.142 \\
\hline\hline
\end{tabular}
\caption{Resultados de detección de anomalías en dataset IOPS. N representa el número de neuronas en las capas B y C para los modelos SNN. Se muestran los mejores resultados obtenidos tras optimización con Optuna para cada configuración.}
\label{tab:resultados-iops-escalabilidad}
\end{table}


\begin{table}[htbp]
\centering
\small
\begin{tabular}{lcccc}
\hline\hline
\textbf{Modelo} & \textbf{Prec} & \textbf{Rec} & \textbf{F1} & \textbf{MSE} \\
\hline
SNN (A--B) & -- & -- & 0.061 & 0.142 \\
SNN (A--B--C) & -- & -- & 0.277 & 0.142 \\
TSFEDL-OhShuLih & 0.280 & 0.984 & 0.436 & 0.081 \\
TSFEDL-KhanZulfiqar & 0.263 & 0.926 & 0.410 & 0.078 \\
TSFEDL-ZhengZhenyu & 0.265 & 0.931 & 0.412 & 0.095 \\
TSFEDL-WeiXiaoyan & 0.244 & 0.858 & 0.380 & 0.142 \\
\hline\hline
\end{tabular}
\caption{Mejores resultados de detección de anomalías en el dataset IOPS. Se muestran las métricas de calidad para cada modelo evaluada.}
\label{tab:resultados-iops}
\end{table}

\subsubsection{CalIt2}
Los resultados en el dataset CalIt2 muestran un rendimiento consistente y aceptable para detección de anomalías (24.80\% de anomalías):
\begin{itemize}
    \item \textbf{Resumen}: La arquitectura híbrida SNN (A--B--C) mantiene un F1-score estable entre 0.417-0.418 en todas las configuraciones de neuronas (100, 200, 400), mostrando robustez ante el escalado. El mejor resultado es F1 = 0.4176 con n\_100, superando significativamente al modelo SNN base (F1 = 0.061) pero quedando por debajo de TSFEDL (F1 = 0.704).
    \item \textbf{Tendencias}: Optuna seleccionó consistentemente kernels \texttt{mexican\_hat} y \texttt{gaussian} con tamaños 5-7, procesamiento \texttt{direct}, y valores de threshold en el rango [-66.8, -67.7]. Los mejores trials se encontraron típicamente entre las evaluaciones 32-75, con duración promedio de 14-42 minutos.
    \item \textbf{Errores}: Recall muy alto (99.4-100\%) pero precisión moderada (26.4\%), resultando en una estrategia conservadora que prioriza la detección de todas las anomalías a costa de generar falsos positivos, apropiada para aplicaciones donde el costo de perder una anomalía es alto.
\end{itemize}

\begin{table}[htbp]
\centering
\small
\begin{tabular}{lccccc}
\hline\hline
\textbf{Modelo} & \textbf{N} & \textbf{Prec} & \textbf{Rec} & \textbf{F1} & \textbf{MSE} \\
\hline
SNN (A--B) & 100 & -- & -- & 0.061 & 0.732 \\
SNN (A--B) & 200 & -- & -- & 0.061 & 0.734 \\
SNN (A--B) & 400 & -- & -- & 0.062 & 0.730 \\
\hline
SNN (A--B--C) & 100 & 0.264 & 0.999 & 0.418 & 0.732 \\
SNN (A--B--C) & 200 & 0.264 & 1.000 & 0.417 & 0.734 \\
SNN (A--B--C) & 400 & 0.264 & 0.994 & 0.417 & 0.730 \\
\hline\hline
\end{tabular}
\caption{Resultados de detección de anomalías en dataset CalIt2. N representa el número de neuronas en las capas B y C para los modelos SNN. Se muestran los mejores resultados obtenidos tras optimización con Optuna para cada configuración.}
\label{tab:resultados-iops-escalabilidad}
\end{table}


\begin{table}[htbp]
\centering
\small
\begin{tabular}{lcccc}
\hline\hline
\textbf{Modelo} & \textbf{Prec} & \textbf{Rec} & \textbf{F1} & \textbf{MSE} \\
\hline
SNN (A--B) & -- & -- & 0.061 & 0.142 \\
SNN (A--B--C) & 0.264 & 0.999 & 0.418 & 0.732 \\
TSFEDL-OhShuLih & 0.543 & 1.000 & 0.704 & 0.036 \\
TSFEDL-KhanZulfiqar & 0.543 & 1.000 & 0.704 & 0.026 \\
TSFEDL-ZhengZhenyu & 0.543 & 1.000 & 0.704 & 0.034 \\
TSFEDL-WeiXiaoyan & 0.522 & 0.970 & 0.679 & 0.090 \\
\hline\hline
\end{tabular}
\caption{Mejores resultados de detección de anomalías en el dataset CalIt2. Se muestran las métricas de calidad para cada modelo evaluada.}
\label{tab:resultados-iops}
\end{table}

\subsection{Comparación con modelos TSFEDL}
Se compararon los siguientes modelos de la librería TSFEDL basados en los resultados de los experimentos de evaluación:
\begin{itemize}
    \item \textbf{OhShuLih}: Modelo de redes neuronales para detección de anomalías en series temporales
    \item \textbf{KhanZulfiqar}: Enfoque basado en deep learning para análisis temporal
    \item \textbf{ZhengZhenyu}: Arquitectura híbrida para detección de patrones anómalos
    \item \textbf{WeiXiaoyan}: Modelo especializado en series temporales multivariadas
\end{itemize}
Observaciones:
\begin{itemize}
    \item \textbf{Rendimiento en IOPS}: Los modelos TSFEDL superan a las SNN con F1-scores en el rango 0.380-0.436 vs 0.061-0.277 de las SNN. OhShuLih obtuvo el mejor resultado (F1=0.436), mientras que la mejor SNN híbrida alcanzó F1=0.277, cerrando parcialmente la brecha respecto al modelo SNN base (F1=0.061).
    \item \textbf{Rendimiento en CalIt2}: Los modelos TSFEDL mantienen ventaja con F1-scores de 0.679-0.704 vs 0.061-0.418 de las SNN. La SNN híbrida (F1=0.418) se acerca más a los baselines que en IOPS, pero los tres mejores modelos TSFEDL (F1=0.704) siguen manteniendo una ventaja de 68\%.
    \item \textbf{Eficiencia}: Las SNN mantienen ventajas potenciales en eficiencia computacional debido a la sparsity inherente de los impulsos discretos, aunque la validación empírica de estas ventajas requiere hardware neuromórfico especializado.
\end{itemize}

\subsection{Importancia de hiperparámetros}
% TODO: Añadir figuras/tablas de Optuna si se dispone
\begin{itemize}
    \item \textbf{Capa convolucional}: impacto de \texttt{kernel\_type}, \texttt{kernel\_size}, \texttt{sigma}.
    \item \textbf{Procesamiento}: comparación entre \texttt{direct}, \texttt{weighted\_sum}, \texttt{max}.
    \item \textbf{Neuronales}: sensibilidad a \texttt{threshold} y \texttt{decay}.
\end{itemize}

\subsection{Eficiencia y escalabilidad}
% TODO: Rellenar con medidas de tiempo reales
\begin{itemize}
    \item \textbf{CPU vs GPU}: \textit{[ganancia observada en entrenamiento e inferencia]}.
    \item \textbf{Tiempo por ventana T=250}: \textit{[ms/ventana]}.
    \item \textbf{Memoria}: \textit{[pico de RAM/VRAM]}.
\end{itemize}

\subsection{Figuras recomendadas}
% TODO: Sustituir rutas por imágenes reales exportadas del pipeline
\begin{figure}[htbp]
    \centering
    \includegraphics[width=0.95\textwidth]{Imagenes/optuna_history_placeholder.png}
    \caption{Histórico de mejores valores por \textit{trial} (Optuna).}
    \label{fig:optuna-history}
\end{figure}

\begin{figure}[htbp]
    \centering
    \includegraphics[width=0.95\textwidth]{Imagenes/confusion_matrix_placeholder.png}
    \caption{Matriz de confusión en \textit{[dataset]} para \textit{[modelo]}.}
    \label{fig:confusion-matrix}
\end{figure}

\begin{figure}[htbp]
    \centering
    \includegraphics[width=0.95\textwidth]{Imagenes/pr_curve_placeholder.png}
    \caption{Curva Precisión-Recall comparativa.}
    \label{fig:pr-curve}
\end{figure}

\subsection{Amenazas a la validez}
\begin{itemize}
    \item \textbf{Particionado}: el uso de split 50/50 temporal puede no capturar todas las variaciones; considerar validación cruzada temporal.
    \item \textbf{Etiquetado}: expansión de etiquetas puede sesgar la estimación de \textit{recall}; justificar en función del objetivo de alertado temprano.
    \item \textbf{Recursos}: número limitado de \textit{trials} en Optuna por restricciones de tiempo/cómputo.
\end{itemize}

\subsection{Resumen}
Los hallazgos clave del análisis experimental son:
\begin{itemize}
    \item \textbf{Rendimiento diferenciado por dataset}: CalIt2 mostró F1-scores estables (0.417-0.418) en la SNN híbrida vs F1≈0.061 en SNN base, mientras que IOPS mostró mejoras substanciales con la capa convolucional (F1=0.277 vs 0.061), evidenciando que el beneficio de la hibridación depende del grado de desbalanceo.
    \item \textbf{Efectividad de la capa convolucional**: La arquitectura híbrida demostró mejoras significativas: 6.8x en CalIt2 (0.418 vs 0.061) y 4.6x en IOPS (0.277 vs 0.061), confirmando el valor de la integración CNN-SNN, especialmente en datasets desbalanceados.
    \item \textbf{Selección óptima de kernels}: Optuna seleccionó consistentemente kernels \texttt{mexican\_hat} y \texttt{gaussian} con procesamiento \texttt{direct}, sugiriendo preferencia por suavizado espacial sobre detección de bordes.
    \item \textbf{Eficiencia computacional**: Las SNN presentan potencial para aplicaciones de tiempo real debido a la sparsity inherente de los impulsos discretos, con ventajas teóricas en eficiencia energética que requieren validación en hardware neuromórfico.
    \item \textbf{Competitividad con TSFEDL}: Los baselines TSFEDL mantienen ventaja sobre las SNN, pero la brecha se redujo significativamente con la capa convolucional: IOPS (F1=0.436 vs 0.277) y CalIt2 (F1=0.704 vs 0.418), demostrando el potencial de las arquitecturas híbridas SNN-CNN.
\end{itemize}